\uniterablesection{Висновки}

Під час передкваліфікаційної практики опрацьовано предметну область інтелектуального
моніторингу особистих звичок і спроєктовано напрями вдосконалення раніше створеної системи
Dailu.
Мети практики досягнуто — обґрунтовано інтеграцію технологій штучного інтелекту для аналізу
поведінки користувача та пристрою Інтернету речей для автоматизованого збору активності з
фізичного світу.

Основні результати практики такі:

\begin{enumerate}[noitemsep, topsep=2pt, leftmargin=*, label=\arabic*)]
    \item проаналізовано предметну область та розглянуто наявні рішення — спеціалізовані трекери звичок, сервіси автоматичного збору активності й застосунки з інтелектуальними функціями; виявлено спільне обмеження: інтелектуальні функції будуються поверх ручного введення даних, а джерела активності не поєднують цифрову й фізичну складові в межах однієї системи;
    \item обґрунтовано доцільність удосконалення системи та сформульовано функціональні й нефункціональні вимоги до неї, зокрема щодо генерування рекомендацій, автоматичної категоризації звичок, діалогового асистента, приймання подій від IoT-пристрою, а також щодо безпеки, розширюваності та приватності;
    \item спроєктовано AI-підсистему як окремий обмежений контекст, що взаємодіє з рештою системи через визначені інтерфейси; за основу обрано використання готової великої мовної моделі з поданням контексту даних користувача за принципом генерації, доповненої вибіркою (RAG), з ізоляцією зовнішнього провайдера, стійкістю до тимчасових збоїв і контролем обсягу переданих даних;
    \item обґрунтовано та спроєктовано IoT-складову — фізичну кнопку-трекер на базі мікроконтролера ESP32, яка надсилає подію натискання на сервер; сервер перетворює її на подію активності й опрацьовує наявним механізмом, спільним з інтеграціями GitHub і Strava, без зміни логіки формування записів;
    \item уточнено технологічний стек удосконаленої системи: серверна частина на платформі .NET (C\#), доступ до даних через Entity Framework Core з провайдером Npgsql, база даних PostgreSQL, клієнт на Angular, інтеграція з LLM-провайдером через API та прошивка IoT-пристрою мовою C++ для ESP32.
\end{enumerate}

Спроєктовані рішення спираються на наявну модульну архітектуру системи й розширюють її двома
новими каналами — інтелектуальним аналізом даних та збором фізичної активності, — не порушуючи
меж відповідальності модулів.
Результати практики становлять основу для реалізації та тестування вдосконаленої системи в
межах кваліфікаційної роботи магістра.
